% CS704 Project 2 -- Deliverable 1: Literature Review
% Limit: 1 page excluding references. Compile with pdflatex (Overleaf: works as-is).
% If the course template differs, paste everything between \begin{document}
% and \end{document} into it.
\documentclass[conference]{IEEEtran}
\usepackage{amsmath}
\usepackage{booktabs}
\usepackage{tabularx}
\usepackage{cite}
\usepackage[hidelinks]{hyperref}

\begin{document}

\title{Step and Orientation Detection with Low-Cost IMUs: A Literature Review}

\author{\IEEEauthorblockN{Inesh Bhanuka}
\IEEEauthorblockA{CS704 Project 2, University of Auckland\\
inesh.bhanuka@gmail.com}}

\maketitle

\section{Introduction}
Pedestrian dead reckoning (PDR) tracks a person from body-worn inertial
sensors without infrastructure. Double-integrating low-cost MEMS
acceleration diverges within seconds unless zero-velocity updates (ZUPT)
are applied, which needs a foot-mounted sensor~\cite{jin2011}. Body-worn
and handheld systems therefore use \emph{step-based} PDR,
$\mathbf{p}_k=\mathbf{p}_{k-1}+L_k[\cos\theta_k,\ \sin\theta_k]^T$,
which splits tracking into the two quantities our device must report:
step events (with length $L_k$) and heading
$\theta_k$~\cite{jin2011,tian2016}.

\section{Step Detection}
All reviewed works find steps as peaks in periodic acceleration. They
differ in the \emph{signal} used and how \emph{thresholds} are set.
Jin~et~al.~\cite{jin2011} rotate acceleration into the East-North-Up
frame using the compass attitude and detect a peak followed by a valley
on the vertical axis. A pair must differ by a calibrated $\Delta$ and be
at least 150\,ms apart. Tian~et~al.~\cite{tian2016} first classify the
carrying mode (holding, swinging, pocket) with a finite-state machine
driven by gyroscope rate and gravity direction. They then low-pass the
mode-specific axis at 5\,Hz and accept peaks above a fixed threshold with
a 0.2\,s minimum spacing and a false-peak rejection window. Lee~et~al.~\cite{lee2015} skip both projection and classification. They use
the orientation-invariant magnitude $|\mathbf{a}|$, with amplitude
thresholds ($\mu\pm\sigma/\alpha$) and timing thresholds that adapt to
recent step statistics. This reached 99.5\% accuracy over seven
placements and 98.7\% when the phone pose changed mid-walk, against
94.4\% for Samsung S-Health. Brajdic and Harle~\cite{brajdic2013}
benchmarked eight step counters on 130 traces (27 subjects, six
placements). Windowed peak detection was best (median error $<$1.3\%),
every method undercounted, and frequency-domain methods degraded in the
back pocket. They recommend gating with a standard-deviation walk
detector, which also covers the stationary false-step case that
Lee~et~al.\ leave open. For step length, the Weinberg model
$L=K\sqrt[4]{a_{\max}-a_{\min}}$ is common~\cite{jin2011,kuang2018}.
Tian's height and step-frequency model gave the lowest spread across
users~\cite{tian2016}.

\section{Orientation Estimation}
Jin~et~al.~\cite{jin2011} used only accelerometer and compass, at under
25\,Hz. Principal component analysis of horizontal acceleration was too
unreliable at that rate, so heading came from integrating horizontal
acceleration over 1.4\,s. Metal cabinets distorting the compass were a
key error source. Fusing two pocketed phones with a MAP constraint cut
the error by up to 73.7\%. Tian~et~al.~\cite{tian2016} averaged
gyro-based yaw over each step (two steps in a pocket, to cancel leg
swing) after unwrapping the $\pm180^\circ$ discontinuity. An
\emph{adaptive offset compensation} step then learns the misalignment
between phone and walking direction. Gyro drift was the dominant error
source. Filter-based fusion handles drift and disturbance more formally.
Yuan~et~al.~\cite{yuan2015} built a quaternion unscented Kalman filter
(UKF): gyro rates propagate the quaternion, and accelerometer
roll/pitch and magnetometer yaw correct it. This avoids gimbal lock and
EKF linearisation error. It achieved $<$10$^\circ$ mean heading error
waist-mounted, running on an STM32F4. Kuang~et~al.~\cite{kuang2018}
used a 15-state EKF with zero-velocity and zero-angular-rate updates in
quasi-static periods. The magnetometer is applied only as a
\emph{relative} heading constraint while the field is detected as
quasi-static, giving RMS heading errors of 2.3$^\circ$ (handheld) to
16.6$^\circ$ (pocket).

\begin{table}[!t]
\centering
\footnotesize
\setlength{\tabcolsep}{2.5pt}
\caption{Comparison of reviewed approaches}
\label{tab:cmp}
\begin{tabularx}{\columnwidth}{@{}l>{\raggedright\arraybackslash}X>{\raggedright\arraybackslash}X>{\raggedright\arraybackslash}p{1.55cm}@{}}
\toprule
Work & Step detection & Orientation & Result \\
\midrule
\cite{jin2011} & Vertical acc.\ (ENU), relative peak--valley & Compass; 2-device MAP fusion & $\le$73.7\% error cut \\
\cite{tian2016} & FSM mode $\rightarrow$ axis; peak + rejection & Gyro yaw, step-avg., offset comp. & 98.91\% pos.\ accuracy \\
\cite{kuang2018} & Acc.\ peak; Weinberg length & EKF, ZARU, quasi-static mag. & 2.3--16.6$^\circ$ RMS \\
\cite{yuan2015} & -- & Quaternion UKF (G+A+M) & $<$10$^\circ$ mean \\
\cite{lee2015} & $|\mathbf{a}|$ peak--valley, adaptive thresholds & -- & 99.5\%; 98.7\% pose change \\
\cite{brajdic2013} & Benchmark; windowed peak best & -- & median err.\ $<$1.3\% \\
\bottomrule
\end{tabularx}
\end{table}

\section{Comparison and Implications}
Table~\ref{tab:cmp} summarises the works and shows three trends. First,
adaptive thresholds on $|\mathbf{a}|$~\cite{lee2015} are robust without
classification. Mode-aware detection~\cite{tian2016} gains accuracy but
needs per-mode tuning and missed 7.5\% of mode transitions. Second,
magnetometer-only heading~\cite{jin2011} is cheap but easily distorted
by nearby steel, a serious risk in a factory. The common remedy is gyro
fusion with magnetic-disturbance gating~\cite{yuan2015,kuang2018}.
Third, placement dominates error: pocket and swinging heading errors are
4--7$\times$ the handheld case~\cite{kuang2018}. Since the SensorTile
will be worn at a fixed body location, we plan two parts. Steps will
come from adaptive peak--valley detection on $|\mathbf{a}|$ with a
variance gate~\cite{lee2015,brajdic2013}. Heading will come from a
tilt-compensated, hard/soft-iron-calibrated magnetometer that is
rejected when the field magnitude or dip angle deviates, optionally
fused with the on-board gyroscope~\cite{yuan2015,kuang2018}.

\begin{thebibliography}{6}
\bibitem{jin2011}
Y.~Jin, H.-S. Toh, W.-S. Soh, and W.-C. Wong, ``A robust dead-reckoning
pedestrian tracking system with low cost sensors,'' in \emph{Proc. IEEE
Int. Conf. Pervasive Comput. Commun. (PerCom)}, Seattle, WA, USA, Mar.
2011, pp.~222--230.

\bibitem{tian2016}
Q.~Tian, Z.~Salcic, K.~I.-K. Wang, and Y.~Pan, ``A multi-mode dead
reckoning system for pedestrian tracking using smartphones,'' \emph{IEEE
Sensors J.}, vol.~16, no.~7, pp.~2079--2093, Apr. 2016.

\bibitem{kuang2018}
J.~Kuang, X.~Niu, and X.~Chen, ``Robust pedestrian dead reckoning based
on MEMS-IMU for smartphones,'' \emph{Sensors}, vol.~18, no.~5, Art.
no.~1391, 2018, doi: 10.3390/s18051391.

\bibitem{yuan2015}
X.~Yuan, S.~Yu, S.~Zhang, G.~Wang, and S.~Liu, ``Quaternion-based
unscented Kalman filter for accurate indoor heading estimation using
wearable multi-sensor system,'' \emph{Sensors}, vol.~15, no.~5,
pp.~10872--10890, 2015, doi: 10.3390/s150510872.

\bibitem{lee2015}
H.-H. Lee, S.~Choi, and M.-J. Lee, ``Step detection robust against the
dynamics of smartphones,'' \emph{Sensors}, vol.~15, no.~10,
pp.~27230--27250, 2015, doi: 10.3390/s151027230.

\bibitem{brajdic2013}
A.~Brajdic and R.~Harle, ``Walk detection and step counting on
unconstrained smartphones,'' in \emph{Proc. ACM Int. Joint Conf.
Pervasive Ubiquitous Comput. (UbiComp)}, Zurich, Switzerland, 2013,
pp.~225--234, doi: 10.1145/2493432.2493449.
\end{thebibliography}

\end{document}
